\section{Gabarito e resoluções comentadas — Raciocínio Lógico}

\begin{solucao}{03.01}\textbf{CERTO.} A única situação que falsifica a condicional é antecedente verdadeiro e consequente falso.\end{solucao}
\begin{solucao}{03.02}\textbf{CERTO.} Essa é a tabela-verdade da implicação material.\end{solucao}
\begin{solucao}{03.03}\textbf{ERRADO.} $q\to p$ é a recíproca e não é equivalente em geral.\end{solucao}
\begin{solucao}{03.04}\textbf{CERTO.} É a contrapositiva.\end{solucao}
\begin{solucao}{03.05}\textbf{ERRADO.} A negação é ``Existe pelo menos um contrato que não foi revisado''.\end{solucao}
\begin{solucao}{03.06}\textbf{CERTO.} Lei do terceiro excluído.\end{solucao}
\begin{solucao}{03.07}\textbf{CERTO.} Uma proposição e sua negação não podem ser simultaneamente verdadeiras.\end{solucao}
\begin{solucao}{03.08}\textbf{CERTO.} Essa é a definição de subconjunto.\end{solucao}
\begin{solucao}{03.09}\textbf{ERRADO.} $1,2\times0,8=0,96$: perda de 4\%.\end{solucao}
\begin{solucao}{03.10}\textbf{CERTO.} Em proporcionalidade inversa, $xy=k$.\end{solucao}
\begin{solucao}{03.11}\textbf{B.} Somente $V\to F$ torna a condicional falsa.\end{solucao}
\begin{solucao}{03.12}\textbf{B.} É a contrapositiva da proposição original.\end{solucao}
\begin{solucao}{03.13}\textbf{CERTO.} De Morgan: $\neg(p\land q)=\neg p\lor\neg q$.\end{solucao}
\begin{solucao}{03.14}\textbf{B.} Negar disjunção produz conjunção das negações.\end{solucao}
\begin{solucao}{03.15}\textbf{B.} Modus ponens.\end{solucao}
\begin{solucao}{03.16}\textbf{CERTO.} Consequente verdadeiro não garante antecedente verdadeiro.\end{solucao}
\begin{solucao}{03.17}\textbf{B.} $\neg\forall x\exists y P=\exists x\forall y\neg P$: existe auditor que não analisou processo algum.\end{solucao}
\begin{solucao}{03.18}\textbf{B.} $70+50-30=90$.\end{solucao}
\begin{solucao}{03.19}\textbf{B.} União $=60+45-25=80$; nenhum $=100-80=20$.\end{solucao}
\begin{solucao}{03.20}\textbf{CERTO.} $0,4\times250=100$.\end{solucao}
\begin{solucao}{03.21}\textbf{C.} $800\times1,15=920$ mil.\end{solucao}
\begin{solucao}{03.22}\textbf{C.} $1,10\times1,20=1,32$: 32\%.\end{solucao}
\begin{solucao}{03.23}\textbf{CERTO.} $0,75\times1,25=0,9375$, ou 6,25\% abaixo.\end{solucao}
\begin{solucao}{03.24}\textbf{C.} Produtividade: $180/(6\times5)=6$ processos por analista-dia. Dez analistas fazem 300 em $300/(10\times6)=5$ dias.\end{solucao}
\begin{solucao}{03.25}\textbf{C.} $x+2x=45\Rightarrow x=15$.\end{solucao}
\begin{solucao}{03.26}\textbf{C.} I e III corretas; II viola De Morgan.\end{solucao}
\begin{solucao}{03.27}\textbf{CERTO.} $\neg(p\to q)=p\land\neg q$.\end{solucao}
\begin{solucao}{03.28}\textbf{B.} Na implicação, antecedente é suficiente e consequente necessário.\end{solucao}
\begin{solucao}{03.29}\textbf{C.} De $q\to r$ e $\neg r$, obtemos $\neg q$; de $p\to q$ e $\neg q$, obtemos $\neg p$.\end{solucao}
\begin{solucao}{03.30}\textbf{CERTO.} O elemento de $B\cap C$ pode estar fora de $A$.\end{solucao}
\begin{solucao}{03.31}\textbf{A.} Auditor $\subseteq$ servidor e terceirizado é disjunto de servidor; logo auditor e terceirizado são disjuntos.\end{solucao}
\begin{solucao}{03.32}\textbf{D.} $120+90+80-50-40-30+20=190$.\end{solucao}
\begin{solucao}{03.33}\textbf{CERTO.} $|A\cap B\setminus C|=50-20=30$.\end{solucao}
\begin{solucao}{03.34}\textbf{C.} $0,8P=800\Rightarrow P=1000$.\end{solucao}
\begin{solucao}{03.35}\textbf{B.} Queda de 5 pontos; relativa a 25: $5/25=20\%$.\end{solucao}
\begin{solucao}{03.36}\textbf{CERTO.} $10/30=1/3$.\end{solucao}
\begin{solucao}{03.37}\textbf{B.} $8\times15=12\times d\Rightarrow d=10$.\end{solucao}
\begin{solucao}{03.38}\textbf{B.} $a+2a=30\Rightarrow a=10$.\end{solucao}
\begin{solucao}{03.39}\textbf{CERTO.} Soma = média $\times$ quantidade = 100.\end{solucao}
\begin{solucao}{03.40}\textbf{C.} $500\times0,60\times0,10=30$.\end{solucao}
\begin{solucao}{03.41}\textbf{CERTO.} Modus tollens: $p\to q$, $\neg q\Rightarrow\neg p$.\end{solucao}
\begin{solucao}{03.42}\textbf{C.} De $p\to q$ e $q$ não se conclui $p$.\end{solucao}
\begin{solucao}{03.43}\textbf{B.} Negar ``existe servidor que revisou todos'' dá ``para todo servidor, existe processo não revisado por ele''.\end{solucao}
\begin{solucao}{03.44}\textbf{CERTO.} Negam-se os quantificadores e o predicado: $\neg\exists\forall=\forall\exists\neg$.\end{solucao}
\begin{solucao}{03.45}\textbf{B.} Só A = 60; só B = 40; exatamente um = 100.\end{solucao}
\begin{solucao}{03.46}\textbf{C.} $0,9\times1,1\times1,05=1,0395$: +3,95\%.\end{solucao}
\begin{solucao}{03.47}\textbf{CERTO.} $20/80=25\%$ e $20/100=20\%$.\end{solucao}
\begin{solucao}{03.48}\textbf{D.} Soma da razão = 10; C recebe $5/10\times400=200$.\end{solucao}
\begin{solucao}{03.49}\textbf{B.} $50+2x=170$ em milhares: $2x=120$, $x=60$.\end{solucao}
\begin{solucao}{03.50}\textbf{B.} ``P somente se A'' significa $P\to A$: aceite é necessário para pagamento.\end{solucao}
